\chapter*{General Introduction}
\addcontentsline{toc}{chapter}{General Introduction}
\markboth{General Introduction}{General Introduction}

\noindent Ideation lies at the start of every software project, engineering initiative, and product design cycle. When teams meet to tackle an open-ended problem, success depends on their ability to generate varied ideas quickly, evaluate trade-offs objectively, and organize raw insights into clear execution steps. For decades, physical whiteboards and adhesive paper notes served as the standard medium for this work. They offered an immediate tactile workspace where participants could sketch diagrams, post notes, and re-order concepts without technical friction. 

However, modern software development and academic collaboration rarely take place entirely within the same room. As hybrid teams and remote development workflows became the norm, organizations turned to digital collaborative boards such as Miro, Mural, and Klaxoon. While these tools successfully replaced physical boards with infinite interactive canvases accessible from web browsers, they introduced new practical bottlenecks that become obvious during intensive sessions:

\begin{itemize}[leftmargin=1.5em]
    \item \textbf{Cognitive Stagnation and Blank-Canvas Friction:} In early phases of a session, teams often encounter cognitive inertia. Junior contributors frequently hesitate to post the first ideas due to evaluation apprehension, while senior members may inadvertently anchor the discussion around a single narrow perspective.
    \item \textbf{Visual Clutter and Information Overload:} Once a session gains momentum and participants post dozens of sticky notes across the board, the canvas quickly descends into visual noise. Related concepts scatter across disparate coordinates, duplicate ideas go unnoticed, and team members spend more time panning and zooming than critically evaluating suggestions.
    \item \textbf{Post-Meeting Synthesis Gap:} The most painful phase of existing brainstorming workflows occurs \textit{after} the meeting ends. Organizing 50 or 100 raw notes into coherent semantic categories, identifying operational risks, comparing conflicting proposals, and formulating concrete action items still requires hours of manual sorting by the facilitator.
    \item \textbf{Monolithic Licensing Constraints:} Existing commercial whiteboards enforce steep per-seat pricing models with coarse access tiers, making them cost-prohibitive for student engineering cohorts, academic laboratories, and small development teams.
\end{itemize}

\noindent To address these practical limitations, we designed and built \textbf{MindSpark}, an AI-powered collaborative brainstorming platform developed as our final Integration Project within the Department of Information Technology at ISET Bizerte. Rather than treating artificial intelligence as a detached chat window or an external add-on, MindSpark integrates Large Language Models directly into the collaborative canvas lifecycle. The platform pairs a low-latency, multi-user visual workspace with a dedicated, decoupled AI microservice capable of executing semantic operations on demand: generating divergent suggestions when teams stall, clustering scattered notes into themed categories, extracting actionable execution plans, evaluating implementation risks, and summarizing board state into structured documents.

Building a production-ready system of this scope required moving beyond typical monolithic student assignments. The platform is engineered as a cloud-native, microservice-oriented distributed architecture. The ecosystem comprises eight Spring Boot services running on Java 21, an Express.js notification engine running on Node.js 24, an Angular 22 single-page application utilizing reactive Signals and native WebSockets, five isolated PostgreSQL database instances following the database-per-service pattern, an Apache Kafka event broker operating in KRaft mode, an enterprise-grade Keycloak 26 identity server, and an observability stack combining Prometheus, Grafana, and Grafana Tempo.

Our project was executed by a six-student engineering team following the Agile Scrum methodology. Over three intensive two-week sprints, the team transformed an initial product backlog of 64 user stories---prioritized according to the MoSCoW method---into a fully containerized platform orchestrated via Docker Compose.

\bigskip
\noindent\textbf{Structure of the Report.} This report documents the analysis, design, architectural decisions, and incremental implementation of MindSpark, organized across six chapters:

\begin{itemize}[leftmargin=1.5em]
    \item \textbf{Chapter 1 (Project Context and Framework):} Outlines the pedagogical framework of the Integration Project at ISET Bizerte, provides a critical study of existing collaborative tools, presents our modeling and Scrum project management methodologies, details the development workstation environment, and establishes the complete technical baseline of our containerized microservice stack.
    \item \textbf{Chapter 2 (State of the Art):} Examines the cognitive psychology underpinning group ideation, reviews real-time web synchronization mechanisms (WebSockets versus CRDTs), evaluates modern LLM capabilities and prompt engineering strategies, and surveys multi-tenant SaaS resource governance models.
    \item \textbf{Chapter 3 (Requirements Analysis and Specification):} Analyzes functional and non-functional requirements, formalizes system actors and global use cases, details core operational scenarios, and introduces the global structural class diagram and domain architecture.
    \item \textbf{Chapter 4 (Sprint 1 --- Working Demo):} Details our first development sprint, covering core user management, invitation-code room creation, sticky note lifecycle, basic canvas interactions, and initial automated test validation.
    \item \textbf{Chapter 5 (Sprint 2 --- Real-Time Collaboration and AI Core):} Explains the second sprint, focusing on the native Spring WebSocket real-time protocol, live cursor broadcasts, the decoupled AI microservice integrating Google Gemini and Groq, concurrency rate limiting, and circuit breaker resilience.
    \item \textbf{Chapter 6 (Sprint 3 --- Full Product, Administration and Quality):} Chronicles our final sprint, introducing the subscription and billing engine integrated with Polar, media uploads, session export history, Kafka-driven asynchronous notifications, distributed tracing via OpenTelemetry, and full test suite execution.
    \item \textbf{General Conclusion:} Synthesizes the engineering achievements of MindSpark, reflects on lessons learned across our distributed architecture, discusses current system boundaries, and proposes concrete avenues for future industrialization.
\end{itemize}
